% 本文件由 LaTeX 排版 Agent 根据有效 translation.json 自动生成。
% author=Gregory of Nyssa; work_id=2906; title=论信德

\chapter{论信德}

% structure: p0001 source=h1 target=omitted reason=page-title-already-rendered-by-parent

致\Person{辛普利丘}。\par

\Person{天主}藉祂的先知命令我们，不可认为任何新神是\Person{天主}，也不可崇拜任何异神。显然，凡不是从永远有的，便称为新；反之，凡不是新的，便称为永远。因此，凡不相信\OriginalTerm{独生天主}{Monogen\={e}s Theos}是从圣父永远有的，并不否认祂是新，因为那不永远有的，显然是新的；而新的不是\Person{天主}，按圣经的话说：\Quote{在你中间不应有任何新神。}因此，他说圣子\Quote{曾经不存在，}就是否认祂的天主性。再者，凡说\Quote{你决不可崇拜异神}的，禁止我们崇拜另一位神；而异神之所以如此称呼，是与我们自己的天主相对而言。那么，我们自己的天主是谁？显然，就是真\Person{天主}。异神是谁？无疑，就是那与真\Person{天主}的性体相异的。因此，如果我们自己的天主是真\Person{天主}，而如异端者所说，独生天主不是出于真\Person{天主}的性体，那么祂就是异神，而不是我们的天主。但福音说，羊\Quote{不会跟随陌生人。}凡说祂是受造的，就使祂与真\Person{天主}的性体相异。那么，那些说祂是受造的，要做什么呢？他们崇拜那受造之物为\Person{天主}，还是不崇拜呢？因为如果他们不崇拜祂，他们就效法犹太人，否认对\Person{基督}的崇拜；如果他们崇拜祂，他们就是拜偶像者，因为他们崇拜的是与真\Person{天主}相异的。但不崇拜圣子与崇拜异神，同样是邪恶的。因此，我们必须说，圣子是真实圣父的真实圣子，这样我们既能崇拜祂，又能避免因崇拜异神而受定罪。但对于那些引用《\WorkTitle{箴言}》中的话：\Quote{上主创造了我\BibleRef{箴 8:28}}，并以为他们由此提出了强有力的论据，证明万物的创造者和造化者被造，我们必须回答说：独生天主为我们\Emph{成了}许多事物。因为祂原是圣言，却成了血肉；祂原是\Person{天主}，却成了人；祂原是无形的，却成了身体；此外，祂还成了\Quote{罪}、\Quote{诅咒}、\Quote{石头}、\Quote{斧头}、\Quote{饼}、\Quote{羔羊}、\Quote{道路}、\Quote{门}、\Quote{磐石}，以及许多类似的事物；祂本性上并非这些中的任何一个，而是为了我们，按照救恩的安排，成了这些。因此，正如祂原是圣言，却为了我们成了血肉；原是\Person{天主}，却成了人；同样，祂原是创造者，却为了我们成了受造物；因为血肉是受造的。因此，正如祂藉先知说：\Quote{上主这样说：那从母胎就塑造了我作祂的仆人\BibleRef{依 49:5}}；同样祂也藉\Person{撒罗满}说：\Quote{上主创造了我，作为祂道路的开始，为祂的化工\BibleRef{箴 8:28}}。因为就如宗徒所说，一切受造物都处于奴役状态。因此，按先知的话，那在童贞女母胎中被塑造的，是仆人，而非主（就是说，按肉体的人，\Person{天主}在其中显示）；同样，在另一处，那被造为祂道路开始的，并不是\Person{天主}，而是那个人，\Person{天主}在其中向我们显示，为更新人类救恩的败坏道路。因此，既然我们在基督内承认两个：一为神性，一为人性（神性按本性，人性按降生成人），我们因此将永远的事归于天主性，将受造的事归于祂的人性。因为如同按先知，祂在母胎中被塑造为仆人，同样按撒罗满，祂藉此奴仆式的受造物显现在肉身中。但当他们说：\Quote{如果祂存在，祂就不是受生的；如果祂是受生的，祂就不存在，}让他们知道，将属于祂肉体起源的属性归于祂的神性本性是不合适的。因为不存在的物体才被产生，\Person{天主}使那些不存在的东西存在，但祂自己并不从不存在中产生。为此，\Person{保禄}也称祂为\Quote{光荣的光辉\BibleRef{希 1:3}}，好使我们知道，正如灯的光与发出光辉的灯同属一性，并且与之结合（因为灯一出现，从它发出的光同时照耀），同样，宗徒在此要我们考虑：圣子出自圣父，圣父从未没有圣子；因为光荣没有光辉是不可能的，正如灯没有光明是不可能的。但显然，正如祂是光辉，证明祂与光荣相关（因为如果光荣不存在，从它发出的光辉也不存在），同样，说光辉\Quote{曾经不存在}就是宣告，当光辉不存在时，光荣也不存在；因为光荣没有光辉是不可能的。因此，就光辉而论，不能说\Quote{如果它存在，它就不是生成的；如果它是生成的，它就不存在，}同样，对圣子说这样的话也是徒劳的，因为圣子就是那光辉。也让那些在圣父与圣子之间谈论\Quote{较小}和\Quote{较大}的人，从\Person{保禄}学习不要衡量不可衡量的事。因为宗徒说，圣子是圣父位格的真像。\BibleRef{希 1:3}那么显然，无论圣父的位格多么大，那位格的真像也同样大；因为真像不可能小于它所显明的位格。伟大的\Person{若望}也这样教导，他说：\Quote{在起初已有圣言，圣言与\Person{天主}同在。}因为他说\Quote{在起初}，而不是\Quote{在起初之后}，表明起初从未没有圣言；他宣称\Quote{圣言与\Person{天主}同在}，意味着圣子相对于圣父毫无欠缺；因为圣言乃是与\Person{天主}的全部存有一起被默观的。因为如果圣言在自己的伟大上有欠缺，以致不能与\Person{天主}的全部存有相关，我们就不得不设想，\Person{天主}超出圣言的那一部分是没有圣言的。但事实上，圣言的整体伟大是与\Person{天主}的整体伟大一起被默观的；因此，在有关天主性的陈述中，谈论\Quote{较大}和\Quote{较小}是不允许的。\par

至於那些說受生者在本性上不像非受生者的人，讓他們從亞當和亞伯的例子學習，不要胡說。因為亞當本人並非按照人類的自然生育而受生；但亞伯是由亞當所生。現在，無疑地，從未受生者被稱為非受生，而藉生育而出者被稱為受生；然而，他未受生的事實並沒有妨礙亞當成為人，亞伯的受生也沒有使他與人的本性有任何不同，而是兩者都是人，儘管一個是藉受生而存在，另一個則未經生育。同樣，在我們關於天主本性的論述中，非受生與受生的事實並不產生本性的差異，而是，正如在亞當和亞伯的例子中，人性是同一個，同樣，在聖父和聖子那裡，天主性也是同一個。\par

如今，關於聖神，褻瀆者也說出與他們論及主時同樣的話，說祂也是受造的。但教會相信，如同關於聖子一樣，關於聖神也同樣相信：祂是非受造的，並且整個受造界藉分享超越自身的善而成為善，而聖神不需要任何事物使祂成為善（因為祂憑本性是善的，正如聖經所見證）；受造界由聖神引導，而聖神給予引導；受造界被治理，而聖神治理；受造界得安慰，而聖神安慰；受造界處於束縛，而聖神賜予自由；受造界被賦予智慧，而聖神賜予智慧的恩寵；受造界分享恩賜，而聖神隨意分施這些恩賜：\BibleQuote{因為這一切都是一個同一的聖神所運行，隨心所欲，個別分配給每一個人。\BibleRef{格前 12:11}} 並且人可以從聖經中找到許多其他證明，指出聖經中應用於聖父和聖子的一切至高無上的、神聖的屬性，也同樣在聖神中被默觀：不朽、有福、良善、智慧、力量、正義、聖潔——每一卓越屬性都被歸於聖神，正如被歸於聖父和聖子一樣，除了那些清楚而明確地將位格彼此區分開來的屬性；我指的是，聖神不被稱為父或子；但聖經中用以稱呼聖父和聖子的所有其他名字，也都同樣應用於聖神。由此，我們理解聖神是超越受造界的。因此，在聖父和聖子被理解的地方，聖神也被理解；因為聖父和聖子是超越受造界的，而我們論證的趨向也為聖神要求這個屬性。所以，一個將聖神置於受造界之上的人，已經接受了正確而健全的道理：因為他將承認我們在聖父、聖子和聖神中所看到的非受造本性是同一的。\par

但是，既然他們根據自己的想法，提出先知的那句話作為聖神受造本性的證據，那句話說：\Quote{祂建立雷聲，創造精神，並向人宣告祂的基督。} 我們必須考慮這一點：先知在建立雷聲時所說的，是創造另一種精神，而非聖神。因為\Quote{雷聲}這個名字，在神秘的語境中被賦予了福音。因此，那些在福音中產生堅定不移的信心的人，從屬肉體變為屬精神，正如主所說的：\BibleQuote{那由肉生的，是肉；那由聖神生的，是精神。} 那麼，是天主藉著建立福音的聲音，使信徒成為精神；而藉著這樣的雷聲從聖神而生、成為精神的人，\Quote{宣告}基督；正如宗徒所說的：\BibleQuote{除非憑聖神，沒有人能說耶穌基督是主。\BibleRef{格前 12:3}}\par

\SourceRecord{Translated by H.A. Wilson.；Nicene and Post-Nicene Fathers, Second Series；Vol. 5.；Edited by Philip Schaff and Henry Wace.；Buffalo, NY: Christian Literature Publishing Co.,；1893.；<http://www.newadvent.org/fathers/2906.htm>.}
